\subsection{斜对称矩阵的普法夫式。}

设 A 是特征 0 的域，\(R = (\alpha_{ij})\) 是 A 上的一个偶数阶 \(2m\) 交错矩阵。以 E 表示向量空间 \(\mathbf{A}^{2m}\) ，\((e_i)\) （ \(i = 1, \ldots, 2m\) ）表示其典范基，\(e\) 表示元素 \(e_1 \wedge e_2 \wedge \ldots \wedge e_{2m}\) （属于 \(\bigwedge^{2m}\) E）。\(m\) 次外幂（这里指双向量 \(u = \sum_{i < j} \alpha_{ij} e_i \wedge e_j \in \bigwedge^2 E\) 的这个外幂）具有 \(\alpha.e\) 的形式，其中 \(\alpha\) 是 A 中一个我们将要计算的元素。元素 \(\bigwedge^m u\) 是形如下式的项之和

\[
\alpha_ {h _ {1} k _ {1}} \alpha_ {h _ {2} k _ {2}} \dots \alpha_ {h _ {m} k _ {m}} e _ {h _ {1}} \wedge e _ {k _ {1}} \wedge e _ {h _ {2}} \wedge e _ {k _ {2}} \wedge \dots \wedge e _ {h _ {m}} \wedge e _ {k _ {m}}\tag{3}
\]

其中 \(h_j < k_j\) （对 \(j = 1, \ldots, m\) ）。这样的项若含有两个相等的 \(e_j\) ，即为零，也就是说，若集合 \(\{h_1, k_1, \ldots, h_m, k_m\}\) 不恰好等于 \(\{1, 2, \ldots, 2m\}\) ，则它为零。此外，若在 (3) 中同时交换 \(e_{h_r}\) 与 \(e_{h_{r+1}}\) 以及 \(e_{k_r}\) 与 \(e_{k_{r+1}}\) ，乘积不变；因此对诸对 \((h_1, k_1), \ldots, (h_m, k_m)\) 施行任何置换，乘积也不变。于是考虑集（而非序列）\(S = \{(h_1, k_1), \ldots, (h_m, k_m)\}\) ，其中的数对 \((h_j, k_j)\) 满足 \(1 \leqslant h_j < k_j \leqslant 2m\) （对 \(j = 1, 2, \ldots, m\) ）；设 \(\mathscr{F}\) 是这些数对所成的集。对 \(S \in \mathscr{F}\) ，令

\[
1 ^ {0}) \varepsilon (S) = 0 \text {   si   } \left\{h _ {1}, k _ {1}, \dots , h _ {m}, k _ {m} \right\} \neq \left\{1, 2, \dots , 2 m \right\};
\]

2°) 在相反的情形，\(\varepsilon(S)=1\) 或 \(\varepsilon(S)=-1\) ，取决于把 \(h_{j}\) 映为 2j-1 且把 \(k_{j}\) 映为 2j（ \(j=1,\ldots,m\) ）的置换是偶置换还是奇置换。

于是前面的注表明 \(\overset{m}{\wedge} u\) 等于

\[
m! \sum_ {S \in \mathfrak {S}} \varepsilon (S) (\prod_ {(h, k) \in S} \alpha_ {h k}) e.\tag{4}
\]

我们现在引入 \(m(2m-1)\) 个未定元 \(X_{hk}\) ，它们以数对 \((h,k)\) （满足 \(1\leqslant h<k\leqslant2m\) ）为指标，并称 P 为 \(\mathbf{Z}\) 上关于诸 \(X_{hk}\) 的由下式定义的多项式

\[
\mathrm{P} ((X _ {h k})) = \sum_ {\mathbf {S} \in \mathfrak {S}} \varepsilon (\mathrm{S}) (\prod_ {(h, k) \in \mathrm{S}} X _ {h k}).\tag{5}
\]

于是我们有

\[
\stackrel {m} {\wedge} u = m! \mathrm{P} ((\alpha_ {h k})). e.\tag{6}
\]

定义 1. --- 给定交错矩阵 \(R = (\alpha_{ij}) (i, j = 1, \ldots, 2m)\) （偶数阶 \(2m\) ，在任意交换环 A 上），称 \(R\) 的普法夫式（记作 \(\mathrm{Pf}(R)\) ）为 A 的元素 \(\mathrm{P}((\alpha_{hk}))\) ，其中 \(1 \leqslant h < k \leqslant 2m\) 。

例. --- 假设：

\[
\alpha_ {1 2} = - \alpha_ {2 1} = \beta_ {1}, \alpha_ {3 4} = - \alpha_ {4 3} = \beta_ {2}, \dots , \alpha_ {2 m - 1, 2 m} = - \alpha_ {2 m, 2 m - 1} = \beta_ {m},
\]

而其余的 \(\alpha_{ij}\) 均为零（参见定理 1）。此时 \(R = (\alpha_{ij})\) 的普法夫式是 \(\beta_1\beta_2\ldots\beta_m\) 。

命题 1. --- 设 R 是交换环 A 上的偶数阶 2m 交错矩阵，P 是 A 上的 2m 阶方阵。则

\[
\operatorname{Pf} (^ {t} P. R. P) = (\det P) \operatorname{Pf} (R).\tag{7}
\]

实际上，先假设 A 是特征 0 的体，并设 \(R = (\alpha_{ij})\) ，\(P = (\beta_{st})\) 。对 \(R\) 联系如下的双向量

\[
u = \sum_ {i <   j} \alpha_ {i j} e _ {i} \wedge e _ {j} = \frac {1}{2} \sum_ {1 \leqslant i, j \leqslant 2 m} \alpha_ {i j} e _ {i} \wedge e _ {j}
\]

\(\wedge\) A \(^{2m}\) 中的，其中 (e \(_i\) ) 表示 A \(^{2m}\) 的典范基；把 \(^t P\) 视为关于基 (e \(_i\) ) 的、自同态 \(f\) （A \(^{2m}\) 的）的矩阵。于是双向量 ( \(\overset{2}{\wedge} f\) )(u) 对应于矩阵 \(^t P.R.P\) ，因为它等于 \(\frac{1}{2}\sum_{i,j,s,t}\beta_{is}\alpha_{ij}\beta_{jt}e_s\wedge e_t\) 。由于扩张 \(\wedge f\) （即 \(f\) 到外代数 \(\wedge\) A \(^{2m}\) 上的扩张）是该代数的自同态（第 III 章，§ 5，第 9 小节），我们有 \(\overset{m}{\wedge}((\overset{2}{\wedge}f)(u)) = (\overset{2m}{\wedge}f)(\overset{m}{\wedge}u)\) ；由于 \(\overset{2m}{\wedge} f\) 是比为 det \(f\) 的位似，于是由 (6) 与定义 1 得出 \(m!\) Pf( \(^t P.R.P\) ) = \(m!\) (det \(P\) )Pf( \(R\) )，从而得到所考虑情形下的 (7)。一般的情形由注意 (7) 的两边都是矩阵 \(R\) 与 \(P\) 的元素的整系数多项式得出（第 IV 章，§ 2，第 5 小节，附注）。

命题 2. --- 对交换环 A 上的每个偶数阶 2m 交错矩阵 R，我们有

\[
\det R = (\operatorname{Pf} (R)) ^ {2}.\tag{8}
\]

实际上，由于 (8) 的两边都是 \(R\) 的元素的整系数多项式，代数恒等式的开拓原理（第 IV 章，§ 2，第 5 小节，附注）表明只需在 A 是特征 0 的体且 det \(R \neq 0\) 的情形下进行证明。若 \(P\) 是 A 上的 \(2m\) 阶可逆方阵，我们有 det( \(^{t}P.R.P\) ) = ( \(\det P\) ) \(^{2}\) ( \(\det R\) ) 以及 \(\mathrm{Pf}(^{t}P.R.P)\) = ( \(\det P\) ) \(\mathrm{Pf}(R)\) （命题 1），因此只需对 \(^{t}P.R.P\) 而不是 \(R\) 证明 (8)。由定理 1 的推论 1，通过适当选取 \(P\) ，可设矩阵 \(^{t}P.R.P\) 有 ( \(\alpha_{ij}\) ) 的形式，其中

\[
\alpha_ {1 2} = - \alpha_ {2 1} = 1, \dots , \alpha_ {2 m - 1, 2 m} = - \alpha_ {2 m, 2 m - 1} = 1,
\]

而其余的 \(\alpha_{ij}\) 均为零（参见例）。现在这个矩阵的行列式等于 1，其普法夫式也等于 1；这就完成了证明。

